% Generated draft: requires mathematical and source review.
\documentclass[11pt]{article}
\usepackage[margin=25mm]{geometry}
\usepackage{fontspec}
\setmainfont{Latin Modern Roman}
\usepackage{amsmath,amssymb,mathtools,unicode-math}
\setmathfont{Latin Modern Math}
\usepackage{enumitem,needspace,xurl,hyperref}
\setlength{\parindent}{0pt}
\setlength{\parskip}{5pt}
\setlength{\emergencystretch}{3em}
\newcommand{\N}{\mathbb{N}}
\newcommand{\Z}{\mathbb{Z}}
\newcommand{\Q}{\mathbb{Q}}
\newcommand{\R}{\mathbb{R}}
\newcommand{\C}{\mathbb{C}}
\newcommand{\F}{\mathbb{F}}
\newcommand{\A}{\mathbb{A}}
\newcommand{\PP}{\mathbb{P}}
\newcommand{\E}{\mathbb{E}}
\newcommand{\eps}{\varepsilon}
\newcommand{\dd}{\,\mathrm{d}}
\begin{document}
\section*{UnsolvedMath ID: 800201}
\textbf{Source problem code:} catalogue entry
\subsection{Definitions and mathematical statement}
\textbf{Unresolved scope.} The supplied excerpt does not state the hypotheses on \(X\) or the coefficient convention for \(CH(X)\). In particular, it does not specify whether \(X\) must be projective, irreducible holomorphic symplectic, or satisfy other conditions from the surrounding source. The formulation below is therefore a statement schema with an explicitly conditional coefficient interpretation.

Under the integral Chow interpretation, write \(CH^*(X)\) for the Chow ring of algebraic cycles modulo rational equivalence with integer coefficients. For an intended smooth complex algebraic variety \(X\), a line bundle \(D\in\operatorname{Pic}(X)\) determines \(c_1(D)\in CH^1(X)\); interpret \(c_i(X)\) as the Chern class \(c_i(T_X)\in CH^i(X)\) of the tangent bundle. Write
\[
\operatorname{cl}_X\colon CH^*(X)\longrightarrow\bigoplus_{r\geq 0}H^{2r}(X,\mathbb Z)
\]
for the integral cycle-class map, and let
\[
R_X=\mathbb Z\bigl[c_1(D),c_i(T_X)\mid D\in\operatorname{Pic}(X),\ i\geq 1\bigr]\subseteq CH^*(X)
\]
be the subring generated by divisor classes and tangent-bundle Chern classes.

\textbf{Conjectural statement schema (Beauville, Conjecture 10).} For every \(X\) satisfying the intended, presently unspecified source hypotheses, every finite list \(D_1,\ldots,D_k\in\operatorname{Pic}(X)\), and every class \(z\in CH^*(X)\) that is a polynomial in \(c_1(D_1),\ldots,c_1(D_k)\) and the classes \(c_i(T_X)\),
\[
\operatorname{cl}_X(z)=0\text{ in }H^*(X,\mathbb Z)
\quad\Longrightarrow\quad
z=0\text{ in }CH^*(X).
\]
Under the integral interpretation above, this says that \(\operatorname{cl}_X|_{R_X}\) is injective. The missing hypotheses on \(X\) and the intended Chow and polynomial coefficients must be confirmed before this schema becomes a fully specified conjecture; the integral cohomology target is retained exactly as supplied.
\subsection{Short English statement}
Beauville's supplied conjecture says that a polynomial in divisor classes and the Chern classes of \(X\) vanishes modulo rational equivalence whenever its integral cohomology class vanishes. The excerpt leaves the permitted varieties \(X\) and the Chow coefficient convention unspecified, so these require review.
\subsection{Sources}
\begin{itemize}
\item Local UnsolvedMath registry: permanent ID \(800201\), exact source problem code \texttt{catalogue entry}, title \emph{Holomorphic Symplectic Geometry --- Conjecture 10}. The supplied record attributes the conjecture to Beauville, with proposed year 2010.
\item Supplied source: Beauville, \emph{Holomorphic Symplectic Geometry: A Problem List} (2010), Conjecture 10, \url{https://arxiv.org/abs/1002.4321}. The record reports extraction method \texttt{arxiv\_latex} and access date 2026-10-16; these are supplied provenance claims, not independently verified access or source verification.
\item The local record labels the problem \texttt{open} and records an update on 2026-10-16. This is supplied registry status, not an independently verified mathematical or literature assessment. Supplied upstream revision: \texttt{d371ee917ac9647fbb06cd6fc529deec6f3dda82}.
\end{itemize}
\end{document}
